\documentclass{article}
\usepackage{graphicx}
\usepackage{amsmath}
\usepackage{amssymb}
\usepackage{geometry}
\geometry{a4paper, margin=1in}

\begin{document}

\title{APM1220\_FA1}
\author{Roland D. Dela Rosa}
\date{August 2026}
\maketitle

\section*{Problem 2.2}
Given the matrices:
\[
A=\begin{pmatrix}-1&3\\ 4&2\end{pmatrix}, \quad B=\begin{pmatrix}4&-3\\ 1&-2\\ -2&0\end{pmatrix}, \quad C=\begin{pmatrix}5\\ -4\\ 2\end{pmatrix}
\]
\begin{itemize}
    \item[(a)] \textbf{5A:}
    \[
    5A = 5 \begin{pmatrix}-1&3\\ 4&2\end{pmatrix} = \begin{pmatrix}-5&15\\ 20&10\end{pmatrix}
    \]
    \item[(b)] \textbf{BA:}
    \[
    BA = \begin{pmatrix}4&-3\\ 1&-2\\ -2&0\end{pmatrix} \begin{pmatrix}-1&3\\ 4&2\end{pmatrix} = \begin{pmatrix}-16&6\\ -9&-1\\ 2&-6\end{pmatrix}
    \]
    \item[(c)] \textbf{A'B':}
    \[
    A'B' = \begin{pmatrix}-1&4\\ 3&2\end{pmatrix} \begin{pmatrix}4&1&-2\\ -3&-2&0\end{pmatrix} = \begin{pmatrix}-16&-9&2\\ 6&-1&-6\end{pmatrix}
    \]
    \item[(d)] \textbf{C'B:}
    \[
    C'B = \begin{pmatrix}5 & -4 & 2\end{pmatrix} \begin{pmatrix}4&-3\\ 1&-2\\ -2&0\end{pmatrix} = \begin{pmatrix}12 & -7\end{pmatrix}
    \]
    \item[(e)] \textbf{Is AB defined?} No. $A$ is $2 \times 2$ and $B$ is $3 \times 2$. The column dimension of $A$ does not match the row dimension of $B$.
\end{itemize}

\section*{Problem 2.4}
\noindent \textbf{Assuming $A^{-1}$ and $B^{-1}$ exist, prove the following properties:}
\begin{itemize}
    \item[(a)] \textbf{Prove that $(A^{\prime})^{-1}=(A^{-1})^{\prime}$:}
    \textit{Proof:} By definition $AA^{-1} = I$. Taking the transpose of both sides gives $(AA^{-1})' = I'$. Using the property $(XY)' = Y'X'$, the left side expands to $(A^{-1})' A'$, and since the transpose of the identity matrix is just $I$, we get:
    \[
    (A^{-1})' A' = I
    \]
    Similarly, starting from $A^{-1}A = I$ and taking the transpose gives $A'(A^{-1})' = I$. Because both products yield the identity matrix, it follows by definition that $(A^{\prime})^{-1} = (A^{-1})^{\prime}$.

    \item[(b)] \textbf{Prove that $(AB)^{-1}=B^{-1}A^{-1}$:}
    \textit{Proof:} To show that $B^{-1}A^{-1}$ is the inverse of $AB$, we multiply them in both orders and verify that they simplify to $I$:
    \[
    (AB)(B^{-1}A^{-1}) = A(BB^{-1})A^{-1} = AIA^{-1} = AA^{-1} = I
    \]
    and
    \[
    (B^{-1}A^{-1})(AB) = B^{-1}(A^{-1}A)B = B^{-1}IB = B^{-1}B = I
    \]
    Since both multiplications result in the identity matrix, the property holds: $(AB)^{-1} = B^{-1}A^{-1}$.
\end{itemize}


\section*{Problem 2.8}
\noindent \textbf{Given the matrix $A = \begin{pmatrix}1&2\\ 2&-2\end{pmatrix}$, find its eigenvalues, normalized eigenvectors, and spectral decomposition:}
\begin{itemize}
    \item \textbf{Eigenvalues:}
    Solve $\det(A - \lambda I) = 0$:
    \[
    \det \begin{pmatrix}1-\lambda & 2\\ 2 & -2-\lambda\end{pmatrix} = \lambda^2 + \lambda - 6 = 0 \implies (\lambda + 3)(\lambda - 2) = 0
    \]
    Thus, $\lambda_1 = 2$ and $\lambda_2 = -3$.

    \item \textbf{Normalized Eigenvectors:}
    \begin{itemize}
        \item For $\lambda_1 = 2$: $(A - 2I)e_1 = 0 \implies -x_1 + 2x_2 = 0$. Choosing $x_2 = 1$ gives $\begin{pmatrix}2\\ 1\end{pmatrix}$, which normalizes to:
        \[
        e_1 = \begin{pmatrix} 2/\sqrt{5} \\ 1/\sqrt{5} \end{pmatrix}
        \]
        \item For $\lambda_2 = -3$: $(A + 3I)e_2 = 0 \implies 2x_1 + x_2 = 0$. Choosing $x_1 = 1$ gives $\begin{pmatrix}1\\ -2\end{pmatrix}$, which normalizes to:
        \[
        e_2 = \begin{pmatrix} 1/\sqrt{5} \\ -2/\sqrt{5} \end{pmatrix}
        \]
    \end{itemize}

    \item \textbf{Spectral Decomposition:}
    Using $A = \sum_{i=1}^2 \lambda_i e_i e_i'$:
    \[
    A = 2 \begin{pmatrix} 2/\sqrt{5} \\ 1/\sqrt{5} \end{pmatrix} \begin{pmatrix} 2/\sqrt{5} & 1/\sqrt{5} \end{pmatrix} 
    + (-3) \begin{pmatrix} 1/\sqrt{5} \\ -2/\sqrt{5} \end{pmatrix} \begin{pmatrix} 1/\sqrt{5} & -2/\sqrt{5} \end{pmatrix}
    \]
    \[
    A = 2 \begin{pmatrix} 4/5 & 2/5 \\ 2/5 & 1/5 \end{pmatrix} - 3 \begin{pmatrix} 1/5 & -2/5 \\ -2/5 & 4/5 \end{pmatrix} = \begin{pmatrix} 1 & 2 \\ 2 & -2 \end{pmatrix}
    \]
\end{itemize}

\section*{Problem 2.11}
Expanding recursively along the first row for a diagonal matrix $A$ of size $p \times p$:
\[
|A| = a_{11}A_{11} = a_{11}a_{22}\cdots a_{pp} = \prod_{i=1}^p a_{ii}
\]

\section*{Problem 2.24}
Let $\Sigma = \begin{pmatrix}4&0&0\\ 0&9&0\\ 0&0&1\end{pmatrix}$:
\begin{itemize}
    \item[(a)] $\Sigma^{-1} = \begin{pmatrix}1/4&0&0\\ 0&1/9&0\\ 0&0&1\end{pmatrix}$
    \item[(b)] Eigenvalues: $\lambda_1 = 4, \lambda_2 = 9, \lambda_3 = 1$ with standard unit eigenvectors $e_1, e_2, e_3$.
    \item[(c)] Eigenvalues of $\Sigma^{-1}$: $1/4, 1/9, 1$ with the same eigenvectors.
\end{itemize}

\section*{Problem 2.30}
\noindent \textbf{Given random vector $X' = [X_1, X_2, X_3, X_4]$ with mean vector $\mu_X' = [4, 3, 2, 1]$ and covariance matrix $\Sigma_X$ partitioned accordingly:}
\begin{itemize}
    \item[(a)] $E(X^{(1)}) = \begin{pmatrix}4\\ 3\end{pmatrix}$
    \item[(b)] $E(AX^{(1)}) = A E(X^{(1)}) = \begin{pmatrix}1 & 2\end{pmatrix}\begin{pmatrix}4\\ 3\end{pmatrix} = 4 + 6 = 10$
    \item[(c)] $Cov(X^{(1)}) = \Sigma_{11} = \begin{pmatrix}3&0\\ 0&1\end{pmatrix}$
    \item[(d)] $Cov(AX^{(1)}) = A \Sigma_{11} A' = \begin{pmatrix}1 & 2\end{pmatrix}\begin{pmatrix}3&0\\ 0&1\end{pmatrix}\begin{pmatrix}1\\ 2\end{pmatrix} = \begin{pmatrix}3 & 2\end{pmatrix}\begin{pmatrix}1\\ 2\end{pmatrix} = 3 + 4 = 7$
    \item[(e)] $E(X^{(2)}) = \begin{pmatrix}2\\ 1\end{pmatrix}$
    \item[(f)] $E(BX^{(2)}) = B E(X^{(2)}) = \begin{pmatrix}1&-2\\ 2&-1\end{pmatrix}\begin{pmatrix}2\\ 1\end{pmatrix} = \begin{pmatrix}(1)(2) + (-2)(1)\\ (2)(2) + (-1)(1)\end{pmatrix} = \begin{pmatrix}0\\ 3\end{pmatrix}$
    \item[(g)] $Cov(X^{(2)}) = \Sigma_{22} = \begin{pmatrix}9&-2\\ -2&4\end{pmatrix}$
    \item[(h)] $Cov(BX^{(2)}) = B \Sigma_{22} B' = \begin{pmatrix}1&-2\\ 2&-1\end{pmatrix}\begin{pmatrix}9&-2\\ -2&4\end{pmatrix}\begin{pmatrix}1&2\\ -2&-1\end{pmatrix} = \begin{pmatrix}33&36\\ 36&48\end{pmatrix}$
    \item[(i)] $Cov(X^{(1)}, X^{(2)}) = \Sigma_{12} = \begin{pmatrix}2&2\\ 1&0\end{pmatrix}$
    \item[(j)] $Cov(AX^{(1)}, BX^{(2)}) = A \Sigma_{12} B' = \begin{pmatrix}1&2\end{pmatrix}\begin{pmatrix}2&2\\ 1&0\end{pmatrix}\begin{pmatrix}1&2\\ -2&-1\end{pmatrix} = \begin{pmatrix}0&6\end{pmatrix}$
\end{itemize}

\section*{Problem 2.41}
\noindent \textbf{Given random vector $X$ with mean vector $\mu_X' = [3, 2, -2, 0]$ and covariance matrix $\Sigma_X = 3I_4$:}
\begin{itemize}
    \item[(a)] \textbf{Find $E(AX)$:}
    \[
    E(AX) = A \mu_X = \begin{pmatrix}1&-1&0&0\\ 1&1&-2&0\\ 1&1&1&-3\end{pmatrix} \begin{pmatrix}3\\ 2\\ -2\\ 0\end{pmatrix} = \begin{pmatrix}1\\ 9\\ 3\end{pmatrix}
    \]

    \item[(b)] \textbf{Find $Cov(AX)$:}
    Using $Cov(AX) = A \Sigma_X A' = 3AA'$:
    \[
    Cov(AX) = 3 \begin{pmatrix}1&-1&0&0\\ 1&1&-2&0\\ 1&1&1&-3\end{pmatrix} \begin{pmatrix}1&1&1\\ -1&1&1\\ 0&-2&1\\ 0&0&-3\end{pmatrix} = \begin{pmatrix}6&0&0\\ 0&18&0\\ 0&0&36\end{pmatrix}
    \]

    \item[(c)] \textbf{Which pairs of linear combinations have zero covariances?}
    All distinct pairs of linear combinations, since all off-diagonal entries in $Cov(AX)$ are zero.
\end{itemize}
\end{document}